% Original analytic teaching diagram, 2026-11-12.
% X=Y~Unif(0,1). Versions chosen: h_I = 1_I for intervals, h_{{a}} = 0 for singletons.
% At a fixed observation y the nested intervals around y all get value 1 while {y} gets 0,
% which no probability measure can satisfy (continuity from above forces mass 1 on {y}).
\definecolor{VB}{HTML}{7188A3}\definecolor{VBF}{HTML}{E5EAF0}
\definecolor{VG}{HTML}{829C8B}\definecolor{VGF}{HTML}{E8EEE9}
\definecolor{VY}{HTML}{B4A16F}\definecolor{VYF}{HTML}{F1ECDD}
\definecolor{VR}{HTML}{A7787D}\definecolor{VRF}{HTML}{F0E5E6}
\definecolor{VP}{HTML}{9588A5}\definecolor{VPF}{HTML}{ECE8F0}
\definecolor{VI}{HTML}{3B4252}\definecolor{VM}{HTML}{6B7280}\definecolor{VL}{HTML}{4C4D4F}
\begin{tikzpicture}[x=1mm,y=1mm,>=Stealth,font=\figurefont\fontsize{8}{10}\selectfont,
lab/.style={text=VM,font=\figurefont\fontsize{7.5}{9}\selectfont}]
\def\L{80}\def\y{44}
% axis
\draw[VL,line width=.7pt] (0,0) -- (\L,0);
\draw[VL,line width=.6pt] (0,-1.2) -- (0,1.2) node[below=2.2mm,text=VI] {$0$};
\draw[VL,line width=.6pt] (\L,-1.2) -- (\L,1.2) node[below=2.2mm,text=VI] {$1$};
\fill[VR] (\y,0) circle (1.1);
\node[text=VI] at (\y,-4.2) {$y$};
% nested intervals
\foreach \i/\w in {1/30,2/18,3/10,4/5}{
  \draw[VB,line width=1.1pt] (\y-\w,\i*4.5) -- (\y+\w,\i*4.5);
  \draw[VB,line width=.6pt] (\y-\w,\i*4.5-1) -- (\y-\w,\i*4.5+1);
  \draw[VB,line width=.6pt] (\y+\w,\i*4.5-1) -- (\y+\w,\i*4.5+1);
  \node[text=VI,anchor=west] at (\y+33,\i*4.5) {$h_{I_\i}(y)=1$};
}
\node[text=VI] at (\y,22.5) {$\vdots$};
% singleton
\draw[VR,line width=1pt] (\y-1.3,28.2) -- (\y+1.3,30.8);
\draw[VR,line width=1pt] (\y-1.3,30.8) -- (\y+1.3,28.2);
\node[text=VI,anchor=west] at (\y+33,29.5) {$h_{\{y\}}(y)=0$};
\draw[VL!60,line width=.4pt,densely dotted] (\y,1.5) -- (\y,27);
\end{tikzpicture}
